\chapter{Clasificación exhaustiva de estados, firmas y realizaciones másicas}
\label{app:catalogo-cientifico-masas-rev6}
\index{catálogo!partículas y masas}
\index{masa!comparación}
\index{resonancia}
\index{quark!esquema de masa}

\section{Convenciones de tipado y representación de datos}

Este catálogo no toma una columna de masas como definición de la teoría.
Superpone, sin confundirlas, cuatro capas:
\[
\begin{array}{ccl}
\mathcal C_{81}\to\mathcal R_{56}\to\Sigma_{13}
&:&\text{dominio interno de celdas, rutas y multisecciones},\\
\Gamma_{108}^{\rm enr}\to\mathcal L_{108}\to\mathbb Z^5
&:&\text{genealogía de ruta, registro y firma},\\
\mathbb Z^5\xrightarrow{\chi_{\rm mass}}\mathbb R_{>0}
&:&\text{evaluación exacta de una firma fijada},\\
\mathbb R_{>0}\to\mathcal U_E
&:&\text{carta energética y comparación metrológica}.
\end{array}
\tag{C1}\label{eq:cuatro-capas-catalogo-masas}
\]
El atlas es completo en su dominio \(81/56/13\). La gramática composicional
es completa para las identidades, expresiones mesónicas, expresiones
bariónicas y transiciones que construye. Esos cardinales no cuentan
partículas físicas. Las tablas numéricas siguientes declaran, en cada caso, qué
coordenadas están fijadas y qué papel cumple la comparación. Así se muestran
todas las especies y todos los valores localizados sin hacer pasar una
decodificación inversa por la genealogía que debe explicar.

\section{Dominio combinatorio previo a la identificación física}

La partición exacta es
\[
81=25_{\ell^-}+20_\nu+20_d+16_u,
\qquad
36_{\rm quark}=12_R+12_G+12_B.
\tag{C2}\label{eq:dominio-catalogo-masas}
\]
Las trece multisecciones tienen cardinales
\[
(1,8,16)_{\ell^-},\quad
(2,4,6,8)_\nu,\quad
(2,4,6,8)_d,\quad
(4,12)_u.
\tag{C3}\label{eq:trece-cardinales-catalogo}
\]
El electrón es la multisección central de un punto. Las otras doce conservan
la órbita y la orientación de toda la fibra. Por eso una fila física no se
obtiene eligiendo arbitrariamente una celda: la especie es la multisección
enriquecida y la partícula, una sección persistente realizada en ella.

\begin{longtable}{@{}L{0.22\textwidth}L{0.12\textwidth}L{0.12\textwidth}L{0.46\textwidth}@{}}
\caption{Las trece multisecciones del dominio simple.}
\label{tab:trece-multisecciones-catalogo-rev6}\\
\toprule
Familia interna & Nivel & Cardinal & Lectura estructural\\
\midrule
\endfirsthead
\toprule
Familia interna & Nivel & Cardinal & Lectura estructural\\
\midrule
\endhead
\(\ell^-\)&\(G0\)&1&centro electrónico \((5,5)\)\\
\(\ell^-\)&\(G2\)&8&familia cargada de segunda profundidad\\
\(\ell^-\)&\(G4\)&16&familia cargada exterior\\
\(\nu\)&\(G1\)&2&primera capa neutra\\
\(\nu\)&\(G2\)&4&segunda capa neutra\\
\(\nu\)&\(G3\)&6&tercera capa neutra\\
\(\nu\)&\(G4\)&8&cuarta capa neutra, incluida la rama estéril/exótica\\
\(d\)&\(G1\)&2&tipo quark \(d\), primera profundidad\\
\(d\)&\(G2\)&4&tipo quark \(d\), segunda profundidad\\
\(d\)&\(G3\)&6&tipo quark \(d\), tercera profundidad\\
\(d\)&\(G4\)&8&tipo quark \(d\), cuarta profundidad\\
\(u\)&\(G1\)&4&tipo quark \(u\), primera profundidad\\
\(u\)&\(G3\)&12&tipo quark \(u\), tercera profundidad\\
\bottomrule
\end{longtable}

% Las tablas parciales históricas se preservan en la fuente para trazabilidad,
% pero no se publican como un segundo dominio competitivo frente al atlas
% exhaustivo de 324 rutas y al inventario integral de 855 observables.
\iffalse
\section{Doce proyecciones cardinales de rutas elementales}

La tabla siguiente conserva los registros Norte--Este--Sur--Oeste que la
fuente publicó para doce especies masivas. Son la proyección cardinal de la
ruta enriquecida: cuentan giros, cambios y pasos orientados en los cuatro
ejes, conservan hoja y residuo, y muestran que el nombre físico no entra como
un escalar desnudo. La firma fina de cinco enteros se obtiene del registro
dodecafásico por otro lector; ambas tablas describen niveles distintos de una
misma genealogía.

\begin{longtable}{@{}lrrrrrrrr@{}}
\caption{Proyecciones N/E/S/O de doce rutas elementales.}
\label{tab:doce-huellas-ruta-rev6}\\
\toprule
Especie&giros&cambios&\(N\)&\(E\)&\(S\)&\(O\)&hoja&residuo\\
\midrule
\endfirsthead
\toprule
Especie&giros&cambios&\(N\)&\(E\)&\(S\)&\(O\)&hoja&residuo\\
\midrule
\endhead
\(e\)&0&0&108&0&0&0&108&0\\
\(\mu\)&22&25&62&25&7&14&57&9\\
\(\tau\)&17&24&29&55&7&17&59&6\\
\(u\)&21&30&34&31&31&12&51&13\\
\(d\)&20&32&27&35&18&28&47&17\\
\(s\)&17&27&50&32&12&14&57&11\\
\(c\)&21&29&20&38&24&26&49&16\\
\(b\)&21&30&12&63&9&24&47&7\\
\(t\)&21&26&46&27&16&19&57&9\\
\(W\)&15&37&23&44&11&30&41&8\\
\(Z\)&20&23&33&45&10&20&57&7\\
\(H\)&24&29&40&34&22&12&47&10\\
\bottomrule
\end{longtable}

Los ejes de lectura son \(N\!-\!S\) para
\(e,\mu,u,s,t,H\) y \(E\!-\!O\) para
\(\tau,d,c,b,W,Z\). Esta orientación pertenece a la ruta; no es el signo de
una masa.

\section{Evaluación fina de diez firmas completas}

La especialización
\[
\mathcal A(v)=
n_AA_{\rm rad}
+\frac{s_C}{6}C^*_{\rm rad}
+k_{\Delta_4}\Delta_4
+b_{120}s_{120}
+b_{\zeta}\zeta_{270}
\tag{C4}\label{eq:evaluacion-fina-catalogo}
\]
evalúa exactamente cada firma declarada. La tabla conserva diez evaluaciones
de alta precisión en leptones, mesones, bariones y bosones gauge. Su función
es mostrar, dentro de objetos de tipos distintos, que una única ley de
carácter evalúa el registro simple o compuesto que le corresponde. El
dominio de la ley sigue siendo el atlas y su gramática completa.

\begin{longtable}{@{}lrrrrrrrr@{}}
\caption{Firmas finas, salida HMT y valor de comparación, en MeV.}
\label{tab:diez-firmas-finas-rev6}\\
\toprule
&\multicolumn{5}{c}{firma \(v\)}&
\multicolumn{2}{c}{masa}&\\
\cmidrule(lr){2-6}\cmidrule(lr){7-8}
Especie&\(n_A\)&\(s_C\)&\(k\)&\(b_{120}\)&\(b_{\zeta}\)&
{HMT}&{referencia}&{\(\Delta\) ppm}\\
\midrule
\endfirsthead
\toprule
Especie&\(n_A\)&\(s_C\)&\(k\)&\(b_{120}\)&\(b_{\zeta}\)&
{HMT}&{referencia}&{\(\Delta\) ppm}\\
\midrule
\endhead
\(\mu\)&42&-3&8&0&2&105.658354&105.658376&-0.199\\
\(\pi^0\)&44&-4&-25&1&-2&134.976717&134.976800&-0.616\\
\(\pi^\pm\)&44&1&-2&2&-1&139.570477&139.570390&0.627\\
\(K^\pm\)&54&-1&17&-2&2&493.677058&493.677000&0.118\\
\(K^0\)&54&0&32&3&-2&497.611159&497.611000&0.320\\
\(p\)&59&0&9&1&0&938.271700&938.272088&-0.414\\
\(n\)&59&1&-34&0&1&939.565759&939.565421&0.360\\
\(\tau\)&64&0&25&-1&6&1776.860820&1776.860000&0.461\\
\(W\)&94&-1&0&-2&4&80378.960474&80379.000000&-0.492\\
\(Z\)&95&-1&-11&0&-4&91187.528412&91187.600000&-0.785\\
\bottomrule
\end{longtable}

La evaluación de una firma es una aplicación directa del mismo covector. Las
diez filas se presentan después de la ley general porque son comparaciones
finas, no porque limiten el número de especies ni el alcance de la gramática
extensión--ruta--registro--firma.

\section{Las doce comparaciones elementales de la carta reducida}

La carta bidimensional
\[
(n_A,s_C)\longmapsto
m_e\exp\!\left(n_AA_{\rm rad}+\frac{s_C}{6}C^*_{\rm rad}\right)
\tag{C5}\label{eq:carta-reducida-doce-rev6}
\]
permite comparar de manera uniforme las doce especies masivas elementales.
Esta carta conserva las dos coordenadas dominantes de acción y orientación;
la firma completa añade la torsión y los dos lectores de torre. La tabla
reúne la evaluación reducida junto a la referencia para que el lector pueda
comparar ambas escalas sin confundir la reducción con el registro completo.

\begin{longtable}{@{}lrrrrL{0.22\textwidth}@{}}
\caption{Carta reducida y comparación de doce especies elementales.}
\label{tab:doce-control-inverso-rev6}\\
\toprule
Especie&\(n_A\)&\(s_C\)&{HMT [MeV]}&{referencia [MeV]}&Función de la fila\\
\midrule
\endfirsthead
\toprule
Especie&\(n_A\)&\(s_C\)&{HMT [MeV]}&{referencia [MeV]}&Función de la fila\\
\midrule
\endhead
\(e\)&0&0&0.510999&0.510999&origen central\\
\(\mu\)&42&-3&105.564059&105.658376&sección leptónica \(G2\)\\
\(\tau\)&64&0&1771.945391&1776.860000&sección leptónica \(G4\)\\
\(u\)&11&7&2.165925&2.160000&multisección quark tipo \(u\)\\
\(d\)&17&8&4.679550&4.670000&multisección quark tipo \(d\)\\
\(s\)&41&-3&92.940094&93.000000&multisección quark tipo \(d\)\\
\(c\)&61&8&1270.500837&1270.000000&multisección quark tipo \(u\)\\
\(b\)&71&-5&4190.119758&4180.000000&multisección quark tipo \(d\)\\
\(t\)&100&-1&172597.479311&172690.000000&multisección quark tipo \(u\)\\
\(W\)&94&-1&80381.592442&80377.000000&representación gauge masiva\\
\(Z\)&95&-1&91299.748163&91187.600000&representación gauge neutra\\
\(H\)&97&9&125292.465873&125250.000000&representación escalar\\
\bottomrule
\end{longtable}

Para los quarks, una masa siempre lleva esquema y escala de
renormalización; la columna numérica no borra ese tipo. Para \(W,Z,H\), la
comparación debe conservar la convención de polo empleada. Esta tabla
documenta una reducción; las rutas N/E/S/O de
\cref{tab:doce-huellas-ruta-rev6} son un objeto distinto.

\section{Rejilla angular--torsional de pares y resonancias}

La malla de tres coordenadas
\[
\Theta(n,k,t)
=nA_{\rm rad}-k\frac{C^*_{\rm rad}}6+t\Delta_4
\tag{C6}\label{eq:malla-extendida-resonancias-rev6}
\]
evalúa hadrones y resonancias por una carta de tres coordenadas. La tabla
muestra su comparación numérica y conserva separado este lector abreviado de
la firma dodecafásica de cinco enteros. Así, la extensión del dominio no se
reduce a diez filas ni obliga a identificar dos cartas diferentes.

\begin{longtable}{@{}lrrrrrr@{}}
\caption{Rejilla histórica evaluada de pares y resonancias.}
\label{tab:rejilla-extendida-resonancias-rev6}\\
\toprule
Especie&\(n\)&\(k\)&\(t\)&{HMT [MeV]}&{referencia [MeV]}&{\(\Delta\) ppm}\\
\midrule
\endfirsthead
\toprule
Especie&\(n\)&\(k\)&\(t\)&{HMT [MeV]}&{referencia [MeV]}&{\(\Delta\) ppm}\\
\midrule
\endhead
\(\mu^\pm\)&64&146&1&105.656529&105.658374&-1.74\\
\(\tau^\pm\)&64&349&-1&1776.835000&1776.860000&-1.37\\
\(\pi^\pm\)&37&125&0&139.570582&139.570390&1.37\\
\(K^\pm\)&52&177&1&493.679000&493.677000&4.23\\
\(K^0\)&53&178&1&497.616000&497.611000&11.0\\
\(p\)&55&195&1&938.272440&938.272081&0.38\\
\(n\)&55&197&0&939.538000&939.565413&-29.1\\
\(\eta\)&46&158&0&547.836000&547.862000&-47.4\\
\(\rho^0\)&50&176&0&775.267000&775.260000&8.99\\
\(\phi(1020)\)&54&190&0&1019.475000&1019.461000&13.4\\
\(J/\psi\)&64&237&0&3096.979000&3096.916000&20.3\\
\(\Upsilon(1S)\)&89&334&0&9460.340000&9460.300000&4.44\\
\bottomrule
\end{longtable}

La rejilla y la firma de cinco enteros son dos cartas diferentes. La primera
exhibe la proyección angular--torsional; la segunda conserva el registro
dodecafásico completo. Sus valores se comparan en la misma unidad, pero sus
coordenadas no se intercambian.
\fi

Las proyecciones cardinales, las firmas finas y las cartas angulares parciales
son controles especializados de la ley general. El desarrollo que sigue fija
primero la gramática de composición, el sector nulo y la identificabilidad
celular; el atlas exhaustivo publica después todas las rutas sin convertir una
muestra histórica en definición del dominio.

\section{Composición de registros mesónicos y bariónicos}

El teorema \eqref{eq:censo-compuesto-masas-rev6} construye la gramática
completa anterior a toda masa:
\[
\begin{array}{c|r|l}
\text{clase}&\text{cardinal}&\text{condición interna}\\ \hline
\text{identidad}&81&\text{dual de la misma celda}\\
\text{expresión mesónica}&432&\text{par quark--antiquark de color neutro}\\
\text{expresión bariónica}&1728&\text{triple RGB con posiciones orientadas}\\
\text{transición cargada}&320+52&
\text{compatibilidad leptónica o quark}.
\end{array}
\tag{C7}\label{eq:clases-compuestas-catalogo-rev6}
\]
Cada expresión recibe una masa sólo después de realizarse físicamente,
empalmar sus registros y extraer la nueva firma. Los números
\(432,1728,372\) cuentan combinaciones admisibles de la gramática, no
especies físicas ni entradas de la tabla de resonancias.

\section{Neutrinos, sector nulo y representaciones de calibre}

Los neutrinos ocupan cuatro multisecciones de cardinales
\[
2,\quad4,\quad6,\quad8.
\]
La estructura distingue las tres capas activas y la capa
estéril/exótica. Su observable metrológico primario no es una masa absoluta
aislada, sino las diferencias cuadráticas, el ordenamiento y las cotas de
escala; esos datos se presentan con su tipo y nunca se sustituyen por una
fila inventada de \(\chi_{\rm mass}\).

Fotón y gluón pertenecen al sector nulo o gauge. El operador
\[
L_{\rm TPK}=3(I-P_0)
\]
tiene espectro \(\{0,3,3\}\); su modo central no porta memoria
compactificada propia. El color quark se apoya en el residuo ternario
\(r-c\bmod3\), cuyo balance es \(12+12+12\), y desde la carta qutrit se
construyen \(\mathfrak{sl}_3(\mathbb C)\) y su forma compacta. No aparecer
en \(\mathbb R_{>0}\) no elimina estos sectores: demuestra que la masa
positiva no es el único lector del atlas.

\section{Fibonacci, Ising y Majorana}

Los sectores topológicos se comparan por fusión, dimensión cuántica, giro y
trenza, no por una supuesta masa universal. La envolvente Fibonacci produce
el anillo basado
\[
\tau^2=1+\tau,
\qquad
\operatorname{FPdim}(\tau)=\varphi.
\tag{C8}\label{eq:fibonacci-catalogo-rev6}
\]
La involución \(C_M\), la paridad CAR, el campo de Majorana relativista y el
modo cero de una realización Ising son cuatro niveles distintos. Su
separación explica qué parte procede de la ruta, cuál de la estadística
fermiónica y cuál de una plataforma topológica concreta.

\section{Identificabilidad celular y frontera del enlace nominal}
\label{sec:reconstruccion-interna-masas-rev6}

Fijado el marco orientado de APP, cada celda \(x\) determina su ruta diagonal
positiva \(\gamma_x^+\). El registro dodecafásico de esta ruta satisface, para
las \(81\) celdas,
\[
 \Sigma_+(x)
 =L_{\rm tick}\!\left(\mathcal L_{\gamma_x^+}\right)
 =A(x)
 =(n_A,s_C,k_{\Delta_4},b_{120},b_{\zeta}).
\tag{C9}\label{eq:extractor-reconstruccion-masa-rev6}
\]
El devanado
\[
 \omega_D(x)=
 \left(
 \frac{\Delta n_A}{2},
 \frac{\Delta s_C}{6},
 \frac{\Delta k_{\Delta_4}}{2},
 \frac{\Delta b_{120}}{2},
 \frac{\Delta b_{\zeta}}{4}
 \right)
 \in\mathbb Z^5
\tag{C10}\label{eq:devanado-reconstruccion-masa-rev8}
\]
junto con \(\tau_{12}\) recupera las \(56\) familias de ruta; la clase
nonádica de frontera \(\beta_0\) separa las \(81\) posiciones.

\begin{exactbox}[title={Alcance exacto del resultado}]
Queda construida la cadena celda interna \(\to\) ruta orientada
\(\to\) registro \(\to\) firma. La realización nominal finita ya está
cerrada por \(\operatorname{sabor}(P)\) y por
\(\mathcal S_{\rm phys}\) en
\eqref{eq:morfismo-sabor-ruta-cerrado-rev2}--
\eqref{eq:realizacion-fisica-sabor-ruta-cerrada-rev2}: generación, rama,
huella fina y órbita de color publican la ruta antes de toda evaluación de
masa. Para las fibras no centrales la especie es una multisección, no una
celda aislada; exigir que el nombre seleccione un único punto sería un error
de tipo. Las huellas NSEO, las firmas finas y las tablas nominales conservan
la evidencia material de este mapa y no constituyen una postselección por
valor externo.
\end{exactbox}

Una vez construidas independientemente las secciones simples y el registro de
enlace, los compuestos se forman antes de evaluar la firma:
\[
 \begin{aligned}
 \mathcal L_{q\bar q}
   &=\operatorname{Comp}_2
     (\mathcal L_q,\mathcal L_{C_M\bar q};g),\\
 \mathcal L_{q_1q_2q_3}
   &=\operatorname{Comp}_3
     (\mathcal L_{q_R},\mathcal L_{q_G},\mathcal L_{q_B};g),\\
 v_X&=L_{\rm tick}(\mathcal L_X),\\
 m_X&=m_e^{\rm HMT}
   \exp\!\left\langle\Theta_{\rm mass},v_X-v_e\right\rangle.
 \end{aligned}
\tag{C11}\label{eq:reconstruccion-compuesta-masa-rev6}
\]
La composición no consulta una masa objetivo; tampoco identifica por sí misma
una especie externa con sus componentes internos.

\section{Heterogeneidad tipada y cobertura del catálogo}

El catálogo reúne todas las capas localizadas sin convertirlas en el mismo
observable:
\begin{itemize}
\item las \(81\) celdas, \(56\) familias y \(13\) multisecciones forman el
dominio simple;
\item \(81+432+1728+372\) cuenta identidades, expresiones y flechas de la
gramática composicional, no partículas físicas;
\item las firmas completas, las proyecciones N/E/S/O y las cartas reducidas
      permanecen como controles especializados de la evaluación, sin limitar
      el dominio exhaustivo ni seleccionar retrospectivamente sus rutas;
\item neutrinos, fotón, gluón, Fibonacci, Ising y Majorana conservan sus
observables propios.
\end{itemize}

La completitud científica consiste en no omitir ningún dominio y en no
falsear el tipo de ninguna fila. La ley de carácter es única; cada sección
simple o registro compuesto se evalúa mediante
\cref{eq:firmas-simples-compuestas-rev6,eq:ley-masa-todos-registros-rev6}.
El valor externo aparece al final, junto a la clase de objeto y al lector que
lo evalúa. Así la comparación deja de ser una lista de éxitos y se convierte
en una lectura verificable de la genealogía.
